\documentclass{resume}

% Redefinir summary para usar titulo en espanol
\renewcommand{\summary}[1]{%
    \sectionTitle{Resumen Profesional}%
    #1
}

% Definir educationItemEs para evitar macros residuales en ingles (Graduating, Relevant coursework)
\newkeycommand{\educationItemEs}[
    university,
    college,
    program,
    graduation,
    grade,
    coursework
]{%
    {%
        \bf \commandkey{program}%
    } \hfill {%
        \commandkey{graduation}%
    } \\
    {%
        \commandkey{university}%
    } \hfill \ifcommandkey{grade}{%
        \commandkey{grade}%
    }{} \\
    {%
        \commandkey{college}%
    } \\
    \ifcommandkey{coursework}{%
        Materias relevantes: \commandkey{coursework}%
    }{}%
}

\begin{document}

% --------- Informacion de Contacto -----------
\introduction[
    fullname={Christian Omar Payán Torróntegui},
    email={omarpayant@gmail.com},
    phone={+52 669 100 3234},
    linkedin={linkedin.com/in/omarpayant/},
    github={github.com/opyntorr}
]

% --------- Resumen Profesional -----------
\summary{Estudiante avanzado de Ingeniería en Robótica y Sistemas Digitales radicado en Mazatlán, especializado en mantenimiento electromecánico de campo, actuadores de potencia y diagnóstico electrónico con multímetro. Experiencia práctica en soporte de terminales y hardware de laboratorio, ensamble de gabinetes industriales con protección ante vibraciones y cableado de potencia. Manejo intermedio de Microsoft Excel (acreditado con curso de 100 horas) para control riguroso de refacciones, bitácoras de servicio y auditoría de inventarios. Disponibilidad para servicio en Mazatlán y rutas del sur de Sinaloa.}

% --------- Educacion ---------
\begin{educationSection}{Educacion}
    \educationItemEs[
        program={Ingeniería en Robótica y Sistemas Digitales (9no Semestre)},
        university={Tecnológico de Monterrey, Campus Guadalajara},
        college={Zapopan, Jalisco},
        graduation={Ago 2022 -- Dic 2026},
        coursework={Máquinas Eléctricas, Mecatrónica, Circuitos de Potencia, Instrumentación y Control}
    ] \vspace{0.2em}

    \educationItemEs[
        program={Curso de 100 Horas en Microsoft Office (Excel Intermedio, Word, PowerPoint)},
        university={Bachillerato General},
        college={Mazatlán, Sinaloa},
        graduation={Certificado Acreditado},
        coursework={Fórmulas, Filtros, Tablas Dinámicas, Control de Inventario de Refacciones y Liquidación}
    ]
\end{educationSection}

% --------- Habilidades Tecnicas -----------
\begin{skillsSection}{Habilidades Tecnicas}
    \skillItem[
        category={Electromecanica y Actuadores},
        skills={Motores DC y Motorreductores, Relevadores, Mecanismos de Espiral y Arrastre, Sensores Ópticos e Inductivos, Cableado y Arneses Industriales, Ensamble Mecánico Resistente a Vibración}
    ] \\
    \skillItem[
        category={Diagnostico y Medicion},
        skills={Multímetro Digital (Voltajes AC/DC, Continuidad, Señales de Control), Fuentes de Alimentación Conmutadas, Detección de Cortocircuitos, Diagnóstico de Tarjetas Electrónicas}
    ] \\
    \skillItem[
        category={Gestion Operativa y Campo},
        skills={Microsoft Excel Intermedio (Fórmulas, Tablas, Bitácoras de Servicio), Control de Inventario de Partes, Movilidad Local y Disponibilidad de Ruta en Sinaloa}
    ]
\end{skillsSection}

% --------- Experiencia Laboral -----------
\begin{experienceSection}{Experiencia Laboral}
    \experienceItem[
        company={Tecnológico de Monterrey — Laboratorio de Manufactura Avanzada},
        location={Zapopan, Jalisco},
        position={Asistente de Laboratorio y Soporte de Hardware},
        duration={Feb 2026 -- Jun 2026}
    ]
    \begin{itemize}
        \vspace{-0.5em}
        \itemsep -9pt {}
        \item Diseñé, fabriqué y ensamblé gabinetes para mini PCs industriales resistentes a vibración continua, organizando el cableado interno para evitar falsos contactos y desgaste mecánico.
        \item Realicé mantenimiento preventivo y correctivo a estaciones de cómputo y terminales de prueba, diagnosticando fallas eléctricas con multímetro y sustituyendo componentes.
        \item Administré el registro de inventario de componentes electrónicos, refacciones y herramientas mediante hojas de cálculo en Excel, manteniendo auditorías con 100\% de precisión.
    \end{itemize}
\end{experienceSection}

% --------- Proyectos Tecnicos -----------
\begin{experienceSection}{Proyectos Tecnicos}
    \projectItem[
        title={Diseño y Diagnóstico de Tarjeta Electrónica DigiMix (con Intel)},
        duration={Ago 2024 -- Dic 2024},
        keyHighlight={Eagle PCB, Microcontroladores STM32/ESP32, Multímetro, Verificación Eléctrica}
    ]
    \begin{itemize}
        \vspace{-0.5em}
        \itemsep -9pt {}
        \item Diseñé en Eagle PCB un circuito embebido con aislamiento de ruido eléctrico, verificando voltajes de alimentación de 3.3V y 5V y continuidad de pistas con multímetro digital.
        \item Validé la comunicación de datos y el funcionamiento estable de sensores y actuadores bajo pruebas continuas de laboratorio, sin fallas lógicas ni sobrecalentamiento.
    \end{itemize}

    \projectItem[
        title={Integración de Sensores y Soportes para Actuadores (Proyecto OMRON)},
        duration={Ago 2025 -- Dic 2025},
        keyHighlight={Sensores Ópticos, Monturas 3D Rígidas, Calibración de Actuadores, Pruebas en Movimiento}
    ]
    \begin{itemize}
        \vspace{-0.5em}
        \itemsep -9pt {}
        \item Diseñé y manufacturé monturas rígidas de precisión para sensores ópticos y actuadores en plataformas móviles, garantizando sujeción firme ante vibraciones de desplazamiento.
        \item Calibré señales de retroalimentación en sensores y actuadores, garantizando lecturas estables y respuesta mecánica sin holgura.
    \end{itemize}
\end{experienceSection}

\end{document}
